\clearpage\SourceLeaf{451}%
\fancyhead[L]{\scriptsize BEMERKUNGEN.}
\pdfbookmark[1]{BEMERKUNGEN.}{bemerkungen-coordinaten}
\begin{center}\textbf{BEMERKUNGEN.}\end{center}

\indent Der Einheit der Koordinaten sowie den verschiedenen Reduktionen der Messung sollten vermutlich die von \textsc{Walbeck} ermittelten Erddimensionen zugrunde gelegt werden.

\textsc{Walbeck} und \textsc{Brummer}. De forma et magnitudine telluris. Aboae 1819 S. 16: ‘Der mittlere Grad oder \(\frac{1}{90}\) eines Meridianquadranten \(=57009^{t}{,}76\). Elliptizität \(=\frac{1}{302{,}78}\)’ [Handschriftliche Bemerkung von \textsc{Gauss}: mittlerer Meridiangrad] \(=\) ‘\(57009^{t}{,}7584\). Der Meter also \(=443^{l}{,}307885\), Verhältnis \(=37299:37300\) Logarithmus \(=0{,}00001164\).’

\textsc{Gauss}. Bestimmung des Breitenunterschiedes zwischen den Sternwarten von Göttingen und Altona. Göttingen 1828. Art. 20. — ‘Wenn man meine Dreiecke so berechnet, als lägen sie auf der Oberfläche eines elliptischen Sphäroids, dessen Dimensionen \textsc{Walbeck} aus der Gesamtheit der bisherigen Gradmessungen abgeleitet hat und das nach unserer besten gegenwärtigen Kenntnis der wirklichen Gestalt \textit{im Ganzen} am nächsten kommt (Abplattung \(\frac{1}{302{,}76}\), der dreihundertsechzigste Teil des Erdmeridians \(=57009{,}746\) Toisen), und dabei von der Polhöhe von Göttingen \(=51^{0}\ 31'\ 37''85\) ausgeht’..

Danach scheint \textsc{Gauss} mehrfach mit der Abplattung \(\frac{1}{302{,}76}\) anstelle der \textsc{Walbeck}schen \(\frac{1}{302{,}78}\) gerechnet zu haben. Tatsächlich liegt die erstere Zahl auch mehreren der noch im handschriftlichen Nachlass vorhandenen Hilfstafeln zugrunde.

\textsc{Gauss} an \textsc{Schumacher}. Göttingen 1830 April 18 ‘Zweite Hilfstafel, Anmerkung: ‘Bei den früher von mir mitgeteilten Koordinaten ist die Einheit \(\frac{1}{10000000}\) des Erdquadranten nach \textsc{Walbeck}s Dimensionen. Um sie also in Koordinaten umzurechnen, denen die Einheit \(\frac{1}{10000000}\) des Erdquadranten nach \textsc{Schmidt}s neuesten zugrunde liegt, muss man sie zuerst mit \(\frac{57009758}{57008551}\) oder mit \(1+\frac{1}{47245}\) multiplizieren.’

\textsc{Gauss}. Bestimmung des Breitenunterschiedes zwischen Göttingen und Altona. Art. 19. ‘Nach der trigonometrischen Verbindung der Sternwarten von Göttingen und Altona liegt Altona 115163,725 Toisen nördlich und 7,211 Toisen westlich von Göttingen. Diese Zahlen beziehen sich auf die Standorte der Meridiankreise. Sie beruhen auf dem Wert der Dreiecksseite Hamburg-Hohenhorn von 13841,815 Toisen, und dieser beruht auf der Basis, die Herr Professor \textsc{Schumacher} 1820 in Holstein gemessen hat. Da der Vergleich der dabei verwendeten Messstangen mit der Normaltoise aber noch nicht \textit{definitiv} abgeschlossen ist, wird die oben angegebene Entfernung künftig noch im gleichen Verhältnis geändert werden müssen wie die Basis selbst. Diese Änderung kann aber jedenfalls nur sehr gering sein.’

Herr Geheimer Staatsrat \textsc{Andrae} in Kopenhagen bemerkt zur Revision der Basis in einem Schreiben vom 5. März 1865, abgedruckt im Generalbericht über die mitteleuropäische Gradmessung für das Jahr 1864, Seite 6. 7. ‘Die von \textsc{Schumacher} angegebene Länge der \textit{Braacker Basis}: 3014,5799 Toisen, die bei den früheren Berechnungen sowohl der dänischen Triangulation als auch der unter der Leitung von \textsc{Gauss} ausgeführten hannoverschen Triangulation verwendet wurde, konnte nur als vorläufiges Ergebnis der Basismessung gelten, da die Reduktion auf den Meeresspiegel und mehrere andere Korrekturen noch nicht berücksichtigt waren. Da diese Reduktionen erheblich größer sind als die Unsicherheiten der mit großer Sorgfalt ausgeführten Messungen selbst, war eine neue Bestimmung nötig. Herr Professor Dr. \textsc{Peters} in Altona hat auch die Güte gehabt, sämtliche Korrekturen ausführlich und mit größter Genauigkeit zu berechnen. Damit ergibt sich die Länge der Basis nun wie folgt:
